% =============================================================
% report.tex - Elaborat k projektu LPI, Tema 9
%   Odporucaci system filmov, serialov, knih a videohier v jazyku Prolog
% =============================================================

\documentclass[11pt,a4paper]{article}

\usepackage[utf8]{inputenc}
\usepackage[T1]{fontenc}
\usepackage[slovak]{babel}
\usepackage[a4paper,margin=2.5cm]{geometry}
\usepackage{listings}
\usepackage{xcolor}
\usepackage{hyperref}
\usepackage{array}
\usepackage{booktabs}
\usepackage{enumitem}

\hypersetup{
  colorlinks=true,
  linkcolor=blue!60!black,
  urlcolor=blue!60!black,
  citecolor=blue!60!black
}

% Listings style pre Prolog
\definecolor{prologkw}{rgb}{0.0,0.0,0.6}
\definecolor{prologcom}{rgb}{0.4,0.5,0.4}
\definecolor{prologstr}{rgb}{0.6,0.0,0.0}

\lstdefinelanguage{prolog}{
  morekeywords={findall,setof,bagof,member,length,sort,format,nl,
                forall,write,writeln,is,not,fail,true,false,
                consult,halt,assert,retract,reverse},
  sensitive=true,
  morecomment=[l]{\%},
  morestring=[b]',
  morestring=[b]"
}
\lstset{
  language=prolog,
  basicstyle=\ttfamily\small,
  keywordstyle=\color{prologkw}\bfseries,
  commentstyle=\color{prologcom}\itshape,
  stringstyle=\color{prologstr},
  showstringspaces=false,
  numbers=left,
  numberstyle=\tiny\color{gray},
  frame=single,
  breaklines=true,
  breakatwhitespace=true,
  tabsize=2,
  captionpos=b
}

\title{
  \large Logické programovanie a inferencia (LPI) \\[0.4em]
  \LARGE Odporúčací systém filmov, seriálov, kníh a videohier v jazyku Prolog \\[0.4em]
  \large Téma 9: Program odporúčajúci filmy/knihy/technológie
}
\author{
  Tím~--- Jozef Belušák, Ján Kapurík, Adam Peško, Robert Kardošš, Branislav Zurian, Dávid Tkáč \\
  Technická univerzita v~Košiciach, Fakulta elektrotechniky a informatiky
}
\date{máj 2026}

\begin{document}
\maketitle

\begin{abstract}
\noindent
V~práci predstavujeme \emph{vysvetliteľný odporúčací systém}
implementovaný v~jazyku SWI-Prolog. Na rozdiel od bežných
filmových odporúčačov rozširujeme doménu na štyri prepojené
typy obsahu --- \textbf{filmy, seriály, knihy a~videohry} --- a~staviame
ich na spoločný dátový model. Vďaka tomu vie systém ponúknuť
\textbf{cross-domain} odporúčania (\uv{páči sa mi film $\Rightarrow$
možno sa mi bude páčiť hra rovnakého žánru}), \textbf{reverzné
dopyty} (\uv{kto všetko by mal rád konkrétnu hru?}) i~tranzitívne
  \textbf{cestovanie po grafe vzťahov} (žáner, krajina, štúdio,
  tvorca) aj constraint výber filmového večera do časového limitu.
  Hlavnou hodnotou je transparentnosť --- pre každé
odporúčanie vie systém vygenerovať vysvetlenie zložené
z~konkrétnych pravidiel.
\end{abstract}

\section{Motivácia problému}

Odporúčacie systémy patria medzi najbežnejšie nasadené aplikácie
umelej inteligencie. Stretávame sa s nimi denne --- Netflix nám
navrhuje filmy, Spotify skladby, Amazon produkty, YouTube ďalšie
video. Komerčné systémy dnes prevažne stavajú na štatistických
modeloch a~neurónových sieťach. Tieto prístupy dosahujú výbornú
presnosť, no majú jednu zásadnú slabinu --- sú
\textbf{netransparentné}. Používateľ dostane odporúčanie, ale
nikto presne nevie, prečo práve toto. Vznikajú otázky o~spravodlivosti,
manipulácii a~\emph{filter bubble} fenoméne.

V~poslednom čase narastá záujem o~\emph{vysvetliteľnú umelú
inteligenciu} (\emph{explainable AI}, XAI). A~práve tu má
\textbf{logické programovanie} prirodzenú výhodu: každé
odporúčanie možno spätne odôvodniť reťazcom použitých pravidiel.
Prolog predstavuje historicky overený nástroj, v~ktorom je
modelovanie pravidlami prvotriednou občianskou kategóriou ---
nie dodatočnou nadstavbou.

Náš projekt ide o~krok ďalej než typické filmové demo. Doménu
sme rozšírili o~\textbf{seriály, knihy a~videohry}, modelované cez jednotný
predikát \texttt{polozka/4}. To otvára celkom novú triedu
dopytov, ktoré klasický filmový odporúčač nevie:

\begin{itemize}[noitemsep]
  \item \emph{Páči sa mi film Dune $\Rightarrow$} odporuč mi
        knihu Dune (filmová adaptácia rovnakého románu, spoločné
        fiktívne univerzum), prípadne hru Mass Effect (zdieľaný
        žáner sci-fi).
  \item \emph{Kto všetko by mal rád hru The Witcher~3?} ---
        reverzný filter ponad pravidlo \texttt{odporucam/2}.
  \item \emph{Z~filmu Dune sa cez graf zdieľaných atribútov
        (žáner $\to$ krajina $\to$ tvorca) dostať k~hre
        Dark~Souls.}
  \item \emph{Vyber mi filmový večer do 300 minút} ---
        optimalizácia nad odporúčanými filmami cez CLP(FD).
\end{itemize}

\section{Aktuálny stav v~oblasti}

Existujú tri hlavné triedy odporúčacích systémov.

\paragraph{Content-based filtering} odporúča položky podobné
tým, ktoré sa používateľovi páčili v~minulosti. Pracuje
s~atribútmi položiek (žáner, autor, kľúčové slová). Výhodou
je nezávislosť od~ostatných používateľov a~dobrá
vysvetliteľnosť. Nevýhodou je obmedzená diverzita --- systém
má sklony zacykliť používateľa do~\emph{filter bubble}.

\paragraph{Collaborative filtering} odporúča to, čo sa páčilo
iným podobným používateľom. Nepotrebuje atribúty položiek, stačia
hodnotenia. Je to prístup, na ktorom je postavený napr. pôvodný
Netflix Prize algoritmus~\cite{aggarwal2016}. Slabou stránkou
je~\emph{cold start problem} --- nového používateľa alebo novú
položku systém nevie kam zaradiť.

\paragraph{Hybridné systémy} kombinujú oba prístupy a~v~praxi
dnes dominujú. Moderné implementácie (YouTube, TikTok, Spotify)
k~tomu pridávajú hlboké neurónové siete, ktoré z~histórie
interakcií učia latentné reprezentácie.

Logické a~\emph{knowledge-based} prístupy tvoria menšinovú,
no dobre etablovanú vetvu~\cite{burke2000,felfernig2008}. Sú
obzvlášť úspešné v~doménach, kde používateľ kupuje položku
zriedka (autá, dovolenky, hypotéky) a~kde je dôležité formálne
odôvodnenie. Náš projekt patrí do~tejto kategórie.

\section{Teoretické východiská}

\subsection{Logické programovanie a~Hornove klauzuly}

Logické programovanie vychádza z~fragmentu predikátovej logiky
prvého rádu --- \textbf{Hornových klauzúl}. Hornova klauzula má
najviac jeden pozitívny literál a~v~Prologu sa zapisuje v~tvare

\begin{lstlisting}[numbers=none]
hlava :- telo1, telo2, ..., teloN.
\end{lstlisting}

\noindent
čo zodpovedá logickej implikácii
$\textit{telo1} \land \dots \land \textit{teloN} \Rightarrow
\textit{hlava}$.

\subsection{SLD-rezolúcia a~unifikácia}

Prolog dokazuje cieľ pomocou \textbf{SLD-rezolúcie}
(\emph{Selective Linear resolution for Definite clauses}).
Pri každom kroku vyberie jeden cieľ z~pracovnej množiny, nájde
klauzulu, ktorej hlavu možno \textbf{unifikovať} s~týmto cieľom,
a~nahradí cieľ telom danej klauzuly~\cite{lloyd1987}.

\subsection{Backtracking a~hľadanie všetkých riešení}

Prolog používa stratégiu \textbf{prehľadávania do~hĺbky
s~backtrackingom}. Ak nejaký podcieľ zlyhá, systém sa vráti
k~poslednému miestu, kde existovali ďalšie alternatívy
(\emph{choice point}), a~vyskúša ich. Vďaka tomu jeden Prolog
dopyt môže priniesť postupne \textbf{všetky} riešenia ---
vlastnosť, ktorú v~iných paradigmách musíme manuálne emulovať
generátormi alebo kolekciami.

\subsection{Negácia ako neúspech (NAF)}

Klasická logika nedovoľuje z~neprítomnosti dôkazu odvodiť
negáciu. Prolog používa pragmatický variant ---
\emph{negation as failure} (\texttt{\textbackslash+}): cieľ
\texttt{\textbackslash+ G} uspeje práve vtedy, keď sa nepodarí
dokázať \texttt{G}. Tento princíp je nutný napr. pre vyjadrenie
\uv{používateľ ešte nekonzumoval položku}:

\begin{lstlisting}[numbers=none]
este_nekonzumoval(P, F) :- \+ uz_konzumoval(P, F).
\end{lstlisting}

\noindent Ide o~tzv. \emph{closed world assumption} ---
predpokladáme, že znalostná báza obsahuje všetko relevantné
a~všetko ostatné je nepravdivé.

\section{Popis riešenia}

\subsection{Architektúra}

Systém má tri logické vrstvy:

\begin{enumerate}[noitemsep]
  \item \textbf{Znalostná báza} --- fakty o~55 položkách
        (15 filmov, 10 seriálov, 15 kníh, 15 hier) a~5~používateľoch.
  \item \textbf{Inferenčné pravidlá} --- pomocné, hlavné
        odporúčanie, skórovanie, vysvetlenie, collaborative
        filtering, cross-domain odporúčanie, filter
        \emph{fanúšikovia}, graf pribuznosti a~constraint
        odporúčanie.
  \item \textbf{Rozhranie} --- demonštračné dopyty
        v~\texttt{examples.pl}.
\end{enumerate}

\subsection{Unifikovaný dátový model \texttt{polozka/4}}
\label{sec:dm}

Kľúčovým návrhovým rozhodnutím je \textbf{zjednotenie štyroch
domén} (film, seriál, kniha, hra) pod jeden predikát:

\begin{lstlisting}[caption={Reprezentácia položiek}]
% polozka(+Id, +Typ, +Nazov, +Rok)
%   Typ in {film, serial, kniha, hra}
polozka(inception, film, 'Inception', 2010).
polozka(breaking_bad, serial, 'Breaking Bad', 2008).
polozka(dune_book, kniha, 'Dune', 1965).
polozka(witcher3,  hra, 'The Witcher 3: Wild Hunt', 2015).
\end{lstlisting}

Spoločné atribúty platia pre všetky typy:

\begin{center}
\begin{tabular}{lll}
\toprule
\textbf{Predikát} & \textbf{Aritnosť} & \textbf{Význam}\\
\midrule
\texttt{zaner/2}       & (Id, Žáner)       & žáner položky\\
\texttt{hodnotenie/2} & (Id, Skóre)       & hodnotenie 1--10\\
\texttt{krajina/2}    & (Id, Krajina)     & krajina pôvodu\\
\texttt{studio/2}     & (Id, Štúdio)      & štúdio / vydavateľstvo / platforma\\
\texttt{tvorca/3}     & (Id, Rola, Meno)  & režisér / showrunner / autor / vývojár\\
\bottomrule
\end{tabular}
\end{center}

Typovo-špecifické sú len rozmery: \texttt{dlzka/2} (minúty
pre filmy), \texttt{epizody/2} (pre seriály), \texttt{strany/2}
(pre knihy) a~\texttt{dlzka\_hry/2} (hodiny pre hry). Tieto sa konzumujú jediným polymorfným
filtrom \texttt{vyhovuje\_rozmeru/3}, ktorý dispatchuje podľa
typu.

\subsection{Hlavné inferenčné pravidlo}

\begin{lstlisting}
odporucam(P, Polozka) :-
    polozka(Polozka, Typ, _, _),
    pouzivatel_chce_typ(P, Typ),
    once(vyhovuje_zaneru(P, Polozka)),
    vyhovuje_hodnoteniu(P, Polozka),
    vyhovuje_rozmeru(P, Polozka, Typ),
    este_nekonzumoval(P, Polozka).
\end{lstlisting}

\noindent
Pravidlo hovorí, že položka je odporúčaná používateľovi vtedy,
keď simultánne platí: typ patrí medzi obľúbené, žáner sa
zhoduje, hodnotenie je dostatočne vysoké, rozmer (dĺžka filmu /
počet strán / dĺžka hry) je v~limite, a~používateľ položku ešte
nekonzumoval. Operátor \texttt{once/1} zabezpečí, že každú
položku vráti backtracking iba raz, aj keď je žánrov spoločných
viacero.

\subsection{Skórovanie a~zoraďovanie}

\begin{lstlisting}
skore_odporucania(P, Polozka, Skore) :-
    hodnotenie(Polozka, H),
    ( vyhovuje_zaneru(P, Polozka) -> BZ = 2 ; BZ = 0 ),
    ( bonus_tvorcu(P, Polozka)   -> BT = 3 ; BT = 0 ),
    ( bonus_krajiny(P, Polozka)  -> BK = 1 ; BK = 0 ),
    ( bonus_studia(P, Polozka)   -> BS = 1 ; BS = 0 ),
    Skore is H + BZ + BT + BK + BS.
\end{lstlisting}

Skóre je súčet hodnotenia ($1$--$10$), bonusu za~žáner ($+2$),
za~obľúbeného tvorcu ($+3$), za~krajinu ($+1$) a~za~štúdio
($+1$). Maximum je teda~$17$.

\subsection{Vysvetľovanie odporúčaní}

Hlavnou pridanou hodnotou logického prístupu je
\textbf{transparentnosť}. Predikát
\texttt{dovody\_odporucania/3} najprv vráti dôvody ako čisté
Prolog termy (typ, žáner, hodnotenie, rozmer, tvorca, krajina
a~štúdio). Predikát \texttt{vysvetli/2} je potom už len
prezentačná vrstva, ktorá tieto dôvody vypíše používateľovi.
Tým sú vysvetlenia zároveň testovateľné v~PlUnit testoch.

\subsection{Collaborative filtering}

Pôvodné jednoduché pravidlo \uv{zdieľajú aspoň dva žánre} sme
rozšírili na skórovanú podobnosť používateľov. Predikát
\texttt{podobnost\_pouzivatelov/4} porovná spoločné žánre, typy
obsahu, tvorcov, krajiny a~štúdiá; žáner a~tvorca majú vyššiu váhu
ako pomocné atribúty. Používatelia sú považovaní za podobných, keď
dosiahnu skóre aspoň~4.

\begin{lstlisting}
podobni_pouzivatelia(P1, P2) :-
    podobnost_pouzivatelov(P1, P2, Skore),
    Skore >= 4.
\end{lstlisting}

\noindent
Samotné \texttt{odporucam\_collaborative/2} už navyše rešpektuje
obľúbený typ cieľovej položky aj typovo špecifický rozmerový limit.

\subsection{Cross-domain odporúčanie}

Hlavná pridaná hodnota oproti pôvodnému filmovému demu je
\emph{prepojenie domén}:

\begin{lstlisting}
cross_domain_odporucam(P, Zdroj, Ciel, Dovod) :-
    uz_konzumoval(P, Zdroj),
    polozka(Zdroj, TypZdroj, _, _),
    polozka(Ciel,  TypCiel,  _, _),
    TypZdroj \= TypCiel,
    pouzivatel_chce_typ(P, TypCiel),
    este_nekonzumoval(P, Ciel),
    vyhovuje_hodnoteniu(P, Ciel),
    vyhovuje_rozmeru(P, Ciel, TypCiel),
    cross_dovod(Zdroj, Ciel, Dovod).

cross_dovod(Z, C, spolocny_zaner(X))   :- zaner(Z,X), zaner(C,X).
cross_dovod(Z, C, spolocny_tvorca(M))  :- tvorca(Z,_,M), tvorca(C,_,M).
cross_dovod(Z, C, spolocna_krajina(K)) :- krajina(Z,K), krajina(C,K).
cross_dovod(Z, C, spolocne_studio(S))  :- studio(Z,S), studio(C,S).
\end{lstlisting}

\noindent
Predikát \texttt{cross\_dovod/3} formálne zachytáva most medzi
dvomi položkami z~rôznych domén --- spoločným žánrom, tvorcom,
krajinou alebo štúdiom. Backtracking automaticky vygeneruje
všetky aplikovateľné mosty.

\subsection{Filter \uv{kto má rád X} (fanúšikovia)}

\begin{lstlisting}
fanusikovia(Polozka, Pouzivatelia) :-
    findall(P,
            ( pouzivatel(P, _),
              ( odporucam(P, Polozka)
              ; uz_konzumoval(P, Polozka) ) ),
            Zoznam),
    sort(Zoznam, Pouzivatelia).
\end{lstlisting}

\noindent
Toto je \emph{reverzný dopyt} nad pravidlom \texttt{odporucam/2}.
Tam, kde štandardný odporúčač pre zadaného používateľa hľadá
položky, my pre zadanú položku hľadáme používateľov.
Rozšírená verzia \texttt{fanusikovia\_s\_dovodmi/2} ku každému
fanúšikovi priloží aj konkrétne dôvody zhody.

\subsection{Cestovanie po grafe vzťahov}
\label{sec:graf}

Položky v~znalostnej báze tvoria graf prepojený cez zdieľané
atribúty. Predikát \texttt{pribuzne\_polozky/3} určuje silu
prepojenia ako počet rôznych zdieľaných atribútov:

\begin{lstlisting}
pribuzne_polozky(P1, P2, Sila) :-
    polozka(P1,_,_,_), polozka(P2,_,_,_), P1 \= P2,
    findall(A, zdielany_atribut(P1, P2, A), AtrZoznam),
    list_to_set(AtrZoznam, Spolocne),
    length(Spolocne, Sila),
    Sila > 0.
\end{lstlisting}

\noindent
Tranzitívna verzia \texttt{retazec/4} vie nájsť cestu medzi
dvomi položkami v~grafe pribuznosti do~zadaného počtu krokov.
Vďaka backtrackingu a~detekcii cyklov (\texttt{\textbackslash+
member(\dots, Acc)}) je hľadanie kompletné a~bezpečné.

\subsection{Adaptácie, univerzá a~hybridné skóre}
\label{sec:hybrid}

Doteraz sme prepájali domény len cez zdieľané \emph{atribúty}
(žáner, krajina, štúdio, tvorca). Tento prístup však necháva
nezachytené dve silné prepojenia: explicitnú \textbf{adaptáciu}
jedného diela na druhé (film podľa knihy, hra podľa knihy)
a~\textbf{spoločné fiktívne univerzum} (Tolkienovo Stredozemie
obsahuje aj \emph{Hobbita}, aj \emph{Pán prsteňov}). Tieto
prepojenia patria do~znalostnej bázy ako samostatné fakty:

\begin{lstlisting}[caption={Adaptácie a~univerzá}]
adaptacia(dune,     dune_book).     % film 2021 podla knihy 1965
adaptacia(witcher3, witcher_book).  % hra podla Sapkowskeho knih

univerzum(dune,         dune_universe).
univerzum(dune_book,    dune_universe).
univerzum(witcher3,     witcher_universe).
univerzum(witcher_book, witcher_universe).
univerzum(hobbit,       middle_earth).
univerzum(lotr,         middle_earth).
\end{lstlisting}

\noindent
Cross-domain odporúčač sa rozšíri o~dva nové \emph{mosty} ---
stačí pridať klauzuly k~existujúcemu predikátu \texttt{cross\_dovod/3}:

\begin{lstlisting}
cross_dovod(Z, C, spolocne_univerzum(U)) :-
    univerzum(Z, U), univerzum(C, U).

cross_dovod(Z, C, adaptacia) :-
    ( adaptacia(Z, C) ; adaptacia(C, Z) ).
\end{lstlisting}

\noindent
\textbf{Hybridné skóre} potom kombinuje content-based skóre
s~ďalšími tromi signálmi:

\begin{lstlisting}
hybridne_skore(P, Polozka, Skore) :-
    skore_odporucania(P, Polozka, Content),
    ( bonus_collab(P, Polozka)    -> BC = 3 ; BC = 0 ),
    ( bonus_adaptacia(P, Polozka) -> BA = 2 ; BA = 0 ),
    ( bonus_univerzum(P, Polozka) -> BU = 2 ; BU = 0 ),
    Skore is Content + BC + BA + BU.
\end{lstlisting}

\noindent
Bonusy modelujú tri intuície: (1)~\emph{collaborative} ---
\uv{niektorý podobný používateľ to už konzumoval}; (2)~\emph{adaptácia}
--- \uv{už si videl iný formát toho istého diela}; (3)~\emph{univerzum}
--- \uv{už si konzumoval iné dielo z~tohto sveta}. Maximum
hybridného skóre je teda~$10 + 2 + 3 + 1 + 1 + 3 + 2 + 2 = 24$.

\subsection{Constraint odporúčanie cez CLP(FD)}

Okrem jednotlivých odporúčaní systém obsahuje aj malý optimalizačný
problém: \texttt{filmovy\_vecer/5}. Zo všetkých filmov odporúčaných
danému používateľovi vyberie neprázdny podzoznam tak, aby sa zmestil
do zadaného časového limitu a~maximalizoval súčet skóre.

\begin{lstlisting}
filmovy_vecer(P, MaxMinut, Vyber, Dlzka, Skore) :-
    ...,
    scalar_product(Dlzky, Bity, #=, Dlzka),
    scalar_product(SkoreFilm, Bity, #=, Skore),
    Dlzka #=< MaxMinut,
    sum(Bity, #>=, 1),
    labeling([max(Skore)], Bity).
\end{lstlisting}

\noindent
Tento príklad ukazuje, že Prolog nemusí len filtrovať fakty, ale
vie riešiť aj kombinatorické úlohy s~tvrdými obmedzeniami.

\section{Ukážky použitia}

Všetky výstupy boli získané spustením
\texttt{swipl -q -g run\_all -t halt examples.pl}.

\subsection{Príklad 1 --- základné odporúčanie}

\begin{lstlisting}[numbers=none]
?- odporucam(jan, P).
  -> [film]  Interstellar (2014)
  -> [film]  Parasite (2019)
  -> [film]  Joker (2019)
  -> [film]  Dune (2021)
  -> [film]  The Shining (1980)
  -> [film]  Arrival (2016)
  -> [film]  Blade Runner (1982)
  -> [kniha] Dune (1965)
  -> [kniha] Foundation (1951)
  -> [kniha] Hyperion (1989)
  -> [kniha] Neuromancer (1984)
  -> [kniha] The Martian (2011)
  -> [kniha] Gone Girl (2012)
  -> [serial] Breaking Bad (2008)
  -> [serial] Stranger Things (2016)
  -> [serial] The Expanse (2015)
  -> [serial] Chernobyl (2019)
  -> [serial] True Detective (2014)
  -> [serial] Westworld (2016)
  -> [serial] The Mandalorian (2019)
\end{lstlisting}

\subsection{Príklad 3 --- vysvetlenie}

\begin{lstlisting}[numbers=none]
?- vysvetli(jan, interstellar).
Polozka: Interstellar (film, 2014)
  - Typ: film (pouzivatel ho chce)
  - Zaner: sci_fi (oblubeny)
  - Hodnotenie: 9/10 (minimum 8)
  - Dlzka filmu: 169 min (maximum 170)
  - reziser: nolan (oblubeny)
  - Krajina: usa (oblubena)
\end{lstlisting}

Toto je presne ten typ výstupu, aký moderné komerčné systémy
len ťažko poskytujú --- každé odporúčanie je odôvodnené
konkrétnymi pravidlami.

\subsection{Príklad 6 --- cross-domain (z~filmu na hru alebo seriál)}

\begin{lstlisting}[numbers=none]
?- cross_domain_odporucam(peter, dune, C, D).
  -> [serial] Arcane             (spolocna_krajina(usa))
  -> [serial] Arcane             (spolocny_zaner(akcia))
  -> [hra] BioShock              (spolocna_krajina(usa))
  -> [hra] BioShock              (spolocny_zaner(akcia))
  -> [hra] BioShock              (spolocny_zaner(sci_fi))
  -> [serial] Black Mirror       (spolocny_zaner(sci_fi))
  -> [serial] Breaking Bad       (spolocna_krajina(usa))
  -> [serial] Chernobyl          (spolocna_krajina(usa))
  -> [hra] Cyberpunk 2077        (spolocny_zaner(sci_fi))
  -> [serial] Game of Thrones    (spolocna_krajina(usa))
  -> [hra] Half-Life 2           (spolocna_krajina(usa))
  -> [hra] Half-Life 2           (spolocny_zaner(akcia))
  -> [hra] Half-Life 2           (spolocny_zaner(sci_fi))
  -> [hra] Horizon Zero Dawn     (spolocny_zaner(akcia))
  -> [hra] Horizon Zero Dawn     (spolocny_zaner(sci_fi))
  -> [hra] The Last of Us        (spolocna_krajina(usa))
  -> [hra] The Last of Us        (spolocny_zaner(akcia))
  -> [serial] The Mandalorian    (spolocna_krajina(usa))
  -> [serial] The Mandalorian    (spolocny_zaner(akcia))
  -> [serial] The Mandalorian    (spolocny_zaner(sci_fi))
  -> [hra] Mass Effect           (spolocny_zaner(sci_fi))
  -> [hra] Portal 2              (spolocna_krajina(usa))
  -> [hra] Portal 2              (spolocny_zaner(sci_fi))
  -> [hra] Red Dead Redemption 2 (spolocna_krajina(usa))
  -> [hra] Red Dead Redemption 2 (spolocny_zaner(akcia))
  -> [serial] The Expanse        (spolocna_krajina(usa))
  -> [serial] The Expanse        (spolocny_zaner(sci_fi))
  -> [serial] True Detective     (spolocna_krajina(usa))
  -> [serial] Westworld          (spolocna_krajina(usa))
  -> [serial] Westworld          (spolocny_zaner(sci_fi))
\end{lstlisting}

\subsection{Príklad 7 --- fanúšikovia hry}

\begin{lstlisting}[numbers=none]
?- fanusikovia(witcher3, U).
U = [maria, tomas].

?- fanusikovia_s_dovodmi(witcher3, P).
P = [ maria-[zaner(fantasy), zaner(rpg),
             tvorca(sapkowski), krajina(polsko),
             studio(cd_projekt_red)],
      tomas-[zaner(akcia), zaner(rpg),
             krajina(polsko), studio(cd_projekt_red)] ].
\end{lstlisting}

\subsection{Príklad 8 --- cesta v~grafe vzťahov}

\begin{lstlisting}[numbers=none]
?- retazec(amelie, dark_souls, 4, Cesta).
  [film]  Amelie  -->
  [film]  Forrest Gump  -->
  [film]  Inception  -->
  [kniha] The Name of the Wind  -->
  [hra]   Dark Souls
\end{lstlisting}

\subsection{Príklad 9 --- adaptácie, univerzá a~hybridné skóre}

Maria už čítala \emph{Hobbita} a~hrala \emph{Witcher~3}. Vďaka
tomu cross-domain odporúčač cez most \emph{adaptacia}
a~\emph{spolocne\_univerzum} navrhne knižného \emph{The Last
Wish}, a~hybridné skóre vyzdvihne aj \emph{Lord of the Rings}
(rovnaké stredozemné univerzum):

\begin{lstlisting}[numbers=none]
?- cross_domain_odporucam(maria, witcher3, C, D),
   ( D = adaptacia ; D = spolocne_univerzum(_) ).
-> [kniha] The Last Wish  (spolocne_univerzum(witcher_universe))
-> [kniha] The Last Wish  (adaptacia)

?- top_hybridne(maria, T).
1. [kniha] The Last Wish      (content=14, hybrid=18)
2. [kniha] The Lord of the Rings (content=16, hybrid=18)
3. [hra] The Legend of Zelda: Breath of the Wild (content=12, hybrid=12)
4. [hra] Disco Elysium        (content=12, hybrid=12)
5. [hra] Cyberpunk 2077       (content=12, hybrid=12)
\end{lstlisting}

\noindent
Príklad ilustruje, ako sa explicitne modelované vzťahy
(\texttt{adaptacia/2}, \texttt{univerzum/2}) prejavia bez
akéhokoľvek štatistického učenia --- iba ako logické dôsledky
nových faktov a~jednej rozšírenej definície \texttt{cross\_dovod}.

\subsection{Príklad 10 --- constraint výber filmového večera}

\begin{lstlisting}[numbers=none]
?- filmovy_vecer(jan, 300, Vyber, Dlzka, Skore).
Dlzka = 286,
Skore = 27,
Vyber = [interstellar, blade_runner].
\end{lstlisting}

\noindent
Výsledok znamená, že z~Janových odporúčaných filmov systém vybral
kombináciu \emph{Interstellar} a~\emph{Blade Runner}: spolu majú 286 minút
a~v~limite 300 minút dosahujú najvyšší súčet skóre.

\section{Diskusia výsledkov}

\paragraph{Silné stránky.} (1)~Transparentnosť a~vysvetliteľnosť ---
systém vie pre každé odporúčanie zdôvodniť, prečo bolo dané.
(2)~Modulárnosť --- pridanie nového pravidla, nového atribútu,
novej domény (napr. hudby) znamená dopísanie pár faktov, bez
zásahu do~existujúcej logiky. (3)~Deklaratívnosť --- kód
popisuje \emph{čo} má platiť, nie \emph{ako} to spočítať.
(4)~Cross-domain prepojenia bezo~zmeny architektúry vďaka
unifikovanému dátovému modelu.

\paragraph{Obmedzenia.} (1)~Malá znalostná báza
(55~položiek, 5~používateľov). V~reálnom Netflixe ide
o~desiatky tisíc položiek a~milióny používateľov.
(2)~Bez učenia z~dát --- pravidlá sú zafixované expertom.
V~hybridnom prístupe by Prolog mohol slúžiť ako symbolický
modul nad ML predikciami. (3)~Collaborative filtering používa
ručne navrhnuté váhy podobnosti namiesto učených latentných
reprezentácií alebo kosínusovej podobnosti nad veľkou maticou
hodnotení.

\paragraph{Možné rozšírenia.}
\begin{enumerate}[noitemsep]
  \item Načítanie znalostnej bázy z~CSV/JSON
        (\texttt{library(csv)}, \texttt{library(http/json)}).
  \item Webové rozhranie cez \texttt{library(http/thread\_httpd)}.
  \item Negatívne vysvetlenia: predikát
        \texttt{preco\_neodporucam/3}, ktorý vráti dôvody zlyhania.
  \item Probabilistické rozšírenie cez ProbLog ---
        kvantifikácia neistoty v~preferenciách.
  \item Pridanie ďalších domén (hudba, podcasty, technológie)
        --- vďaka unifikovanému modelu len ako nové fakty.
\end{enumerate}

\section{Rozdelenie práce v~tíme}

\begin{center}
\begin{tabular}{>{\bfseries}lp{6cm}l}
\toprule
Člen & Zodpovednosť & Súbory\\
\midrule
Člen 1 & Dátový model \texttt{polozka/4}, znalostná báza
         (55~položiek, 5~používateľov) & \texttt{program.pl} (§1)\\
Člen 2 & Hlavné inferenčné pravidlo \texttt{odporucam/2},
         predikáty \texttt{vyhovuje\_*}, polymorfný filter
         rozmerov & \texttt{program.pl} (§2.A--B)\\
Člen 3 & Pokročilá inferencia: skórovanie, top, vysvetlenie,
         collaborative, \textbf{cross-domain},
         \textbf{fanúšikovia}, \textbf{graf vzťahov},
         \textbf{adaptácie / univerzá}, \textbf{hybridné skóre},
         \textbf{constraint odporúčanie} &
         \texttt{program.pl} (§2.C--G)\\
Člen 4 & Demonštračné príklady (10~scenárov), PlUnit regresné
         testy & \texttt{examples.pl}, \texttt{tests.pl},
         \texttt{README.md}\\
Člen 5 & Elaborát, prezentácia, koordinácia tímu &
         \texttt{src/report.tex}\\
\bottomrule
\end{tabular}
\end{center}

Práca prebiehala v~týždenných iteráciách s~krížovou revíziou kódu.

\bibliographystyle{plain}
\bibliography{bibliography}

\end{document}
